\documentclass[12pt, a4paper]{article}

\usepackage{graphicx}
\graphicspath{ {./img/} }
\usepackage[hidelinks]{hyperref}
\usepackage{amsmath}

\newcommand{\projectname}{Tutorial Time (TT)}

\setlength\parindent{0pt}

\usepackage{fancyhdr}
\setlength{\headheight}{15.2pt}
\pagestyle{fancy}
\fancyhf{}
\renewcommand{\headrulewidth}{0px}

\fancyhead[L]{\projectname}
\fancyfoot[R]{\thepage}

\begin{document}

\pagenumbering{arabic}
\fancyfoot[R]{Page \thepage}

\section{Introduction}
This document describes a solution for the problem of tutorial scheduling. This algorithm can also be used for various other scheduling problems (e.g. lab scheduling), or for the set covering problem. \\
The tutorial scheduling problem is as follows: there is a course of many students, each of whom has their own timetable. The lecturer wishes to schedule the minimal number of tutorial slots such that every student in the course can attend at least one of the times.
 Simultaneously, there might be external constraints such as room size (which limits the maxmimal attendance in any singuar tutorial) or TA availability which must be considered. 

\section{Data Representation}
The data consists of a collection of timetables (of both students and TAs) and of time slots. The student timetables are named $s_1, s_2$, etc. and the time slots are named $t_1, t_2$, etc. \\ 
The timtables are represented as binary matrices arranged in such a way that each entry represents a possible timeslot. The rows represent different hours (e.g. 9-10am, 10-11am, etc) and the columns represent days (e.g. Monday, Tuesday, etc). Entries of 1 represent times that person is free, and entries of 0 mean they are busy.\\
Example matrix (5 columns as 5 days of the week, 8 rows as the 8 possible hours):\\

$\begin{bmatrix}
1 & 0 & 0 & 0 & 0\\
0 & 1 & 1 & 0 & 0\\
0 & 0 & 0 & 0 & 0\\
1 & 1 & 1 & 1 & 1\\
1 & 0 & 0 & 0 & 0\\
0 & 1 & 1 & 0 & 0\\
0 & 0 & 0 & 0 & 0\\
1 & 1 & 1 & 1 & 1
\end{bmatrix}$ 

\par
All timetables are then summed together into one matrix in which each entry represents how many students are available at that time. This matrix will be referred to as "the complete matrix" or M in the rest of this document. \\
This matrix is turned into an undirected graph where each node represents a timeslot with a weight of the number of students available at that time. The edges connecting two nodes represent all students which are busy at both times, and that number is the weight assigned to that node. This graph will be called "the graph" or G for the rest of this document.\\

Individual timeslots are also represented as sets of all people who are available at that time, e.g. \[t_1 = (s_1, s_2, s_3)\]
The size of these sets (i.e. the number of students available at a time $t_i$) will be noted as the size of $t_i$, or $|t_i|$\\

Other relevant data points are the total number of students N and the total number of time slots n (in this case n=40).

\section{Base Algorithm}
The first step of the algorithm is to sort all times in decreasing order by their size (when their size is equal, they should be ordered increasingly by the degree of their respective graph node). \\

Then there are two base cases, after which the algorithm iterates over a set of steps until a solution is reached.\\

Base case 1: There are one or more times when everyone is available. This can be noticed by any entry in the complete matrix being equal to the total number of participants or $|t_i|=N$. If this is true, than those times are trivally the best times and can be immediately returned.\\

Base case 2: There are $n \geq 2$ times for which each student is busy in at most one time (e.g. there are 3 times \textrightarrow{} all students are available for 2 or 3 of those times). This will most frequently occur for n=2.
 In the graph, these are 2 or more nodes which have no edges between them. If represented as a colouring problem, they would be the same colour. If this is true, then those times are the optimal times and will be returned. \\  

When neither of these cases occur, the following steps are repeated until all times have been considered: \\
\begin{enumerate}
\item Observe the first (i.e. smallest) time in the queue $t_i$.
\item Create U, the union of all times other than the current $t_i$.
\item If $t_i \subset U$, delete the time $t_i$. 
\end{enumerate}
This is because if a time t is a subset of all other times, all the students free at t are also free at some other time. Due to the times being sorted smallest to largest, this algorithm will preserve the times with the most students available.

\section{Consideration of Constraints}
Oftentimes, external constraints interfere with this algorithm. The constraints considered here are: certain times must be blocked off (e.g. due to TA availability), there is a hard limit on tutorial number (e.g. due to funding), and there is a size limit on tutorials (e.g. due to room size).\\

Limits on tutorial number should hopefully be followed by finding the minimal number of tutorial times, but if that is still too large, the x largest tutorial times can be returned (or all times are returned and the lecturer is warned that they must choose x times to fit into their limit).\\

Times can be blocked off by deleting those sets and entries in schedules before the algorithm executes.\\

Often there is also a minimum number m of TAs that must be present at each tutorial. If there are fixed TAs, their schedules can be considered first and then bad times can be blocked off. More precisely, all TA schedules are added together and then all times where less than than m TAs are available are blocked off as described above. \\

The size limit is more interesting. The solution proposed here features extra data pre-processing. Before the algorithm is run, all times whose size exceeds the tutorial size limit x are pulled aside. \\
For each of these times, the following steps are executed: \\
\begin{enumerate} 
\item Create U, the union of all times other than the current $t_i$.
\item Take the intersection $I = U \cap t_i$.
\item Remove elements from $t_i$ that are in I until $|t_i|=x$.
\end{enumerate}

After this, the algorithm is followed regularly.

\end{document}
